%Daten zur Institution

%\input{Dozentdaten}

%\renewcommand{\fachbereich}{Fachbereich}

%\renewcommand{\dozent}{Prof. Dr. . }

%Klausurdaten

\renewcommand{\klausurgebiet}{ }

\renewcommand{\klausurtyp}{ }

\renewcommand{\klausurdatum}{ . 20}

\klausurvorspann {\fachbereich} {\klausurdatum} {\dozent} {\klausurgebiet} {\klausurtyp}

%Daten für folgende Punktetabelle

\renewcommand{\aeins}{  3  }

\renewcommand{\azwei}{  3   }

\renewcommand{\adrei}{  6  }

\renewcommand{\avier}{  3  }

\renewcommand{\afuenf}{  3  }

\renewcommand{\asechs}{  5  }

\renewcommand{\asieben}{  13  }

\renewcommand{\aacht}{  4  }

\renewcommand{\aneun}{  3  }

\renewcommand{\azehn}{  7  }

\renewcommand{\aelf}{  6  }

\renewcommand{\azwoelf}{  3  }

\renewcommand{\adreizehn}{  6   }

\renewcommand{\avierzehn}{  65  }

\renewcommand{\afuenfzehn}{    }

\renewcommand{\asechzehn}{    }

\renewcommand{\asiebzehn}{    }

\renewcommand{\aachtzehn}{    }

\renewcommand{\aneunzehn}{    }

\renewcommand{\azwanzig}{    }

\renewcommand{\aeinundzwanzig}{    }

\renewcommand{\azweiundzwanzig}{    }

\renewcommand{\adreiundzwanzig}{    }

\renewcommand{\avierundzwanzig}{    }

\renewcommand{\afuenfundzwanzig}{    }

\renewcommand{\asechsundzwanzig}{    }

\punktetabelledreizehn

\klausurnote

\newpage

\setcounter{section}{0}

__NOEDITSECTION__

\inputaufgabeklausurloesung
{3}

{

Definiere die folgenden
\zusatzklammer {kursiv gedruckten} {} {} Begriffe.
\aufzaehlungsechs{Ein
\stichwort {Hausdorff-Raum} {}
$X$.

}{Eine
\stichwort {differenzierbare Abbildung} {}
\maabbdisp {\varphi} { L } { M
} {}
zwischen zwei differenzierbaren Mannigfaltigkeiten
\mathkor {} {L} {und} {M} {.}

}{Das
\stichwort {Tangentialbündel} {}
zu einer differenzierbaren Mannigfaltigkeit $M$.

}{Die
\stichwort {zweite Fundamentalmatrix} {}
zu einer lokalen Parametrisierung
\maabbdisp {} {V} {U \subseteq Y
} {}
\zusatzklammer {mit

\mavergleichskette
{\vergleichskette
{ V
}
{ \subseteq }{ \R^{n-1}
}
{  }{
}
{    }{
}
{   }{
}
}
{}{}{}
offen} {} {}
einer differenzierbaren Hyperfläche

\mavergleichskette
{\vergleichskette
{ Y
}
{ \subseteq }{ \R^n
}
{  }{
}
{    }{
}
{   }{
}
}
{}{}{.}

}{Die \stichwort {äußere Ableitung} {} zu einer stetig differenzierbaren Differentialform
\mathl{\omega \in
{ \mathcal E }^{ k } (  M   )}{} auf einer differenzierbaren Mannigfaltigkeit $M$.

}{Eine \stichwort {der Überdeckung untergeordnete Partition der Eins} {,} wobei

\mavergleichskette
{\vergleichskette
{  X
}
{ = }{ \bigcup_{i \in I} W_i
}
{  }{
}
{    }{
}
{   }{
}
}
{}{}{}
eine
\definitionsverweis {offene Überdeckung}{}{}
eines
\definitionsverweis {topologischen Raumes}{}{}
$X$ ist.
}

}
{

\aufzaehlungsechs{Der
\definitionsverweis {topologische Raum}{}{}
$X$ heißt hausdorffsch, wenn es zu je zwei verschiedenen Punkten
\mathl{x, y \in X}{} zwei
\definitionsverweis {offene Mengen}{}{}
\mathkor {} {U} {und} {V} {}
gibt mit
\mathl{x \in U, \, y \in V}{} und
\mathl{U \cap V = \emptyset}{.}
}{Es seien
\mathkor {} {(U_i,U_i', \alpha_i, i \in I)} {und} {(V_j,V_j', \beta_j, j \in J)} {}
die
\definitionsverweis {Atlanten}{}{}
von
\mathkor {} {M} {und} {N} {.}
Die Abbildung
\maabbdisp {\varphi} { L } { M
} {}
heißt differenzierbar, wenn sie stetig ist und wenn für alle
\mathkor {} {i \in I} {und alle} {j \in J} {}
die Abbildungen
\maabbdisp {\beta_j \circ  \varphi \circ (\alpha_i)^{-1}} {\alpha_i( \varphi^{-1}(V_j) \cap U_i) } { V_j'
} {}
\definitionsverweis {stetig differenzierbar}{}{}
sind.
}{Man nennt die Menge

\mavergleichskettedisp
{\vergleichskette
{TM
}
{ =} {\biguplus_{P \in M} T_PM
}
{ } {
}
{ } {
}
{ } {
}
}
{}{}{,}
versehen mit der Projektionsabbildung
\maabbeledisp {\pi} {TM} {M
} {(P,v)} {P
} {,}
das Tangentialbündel von $M$.
}{Es sei $Y$ durch ein
\definitionsverweis {Einheitsnormalenfeld}{}{}
$N$
\definitionsverweis {orientiert}{}{,}
wobei wir $N$ als Feld auf $V$ auffassen. Dann setzt man

\mavergleichskettedisp
{\vergleichskette
{ h_{i j}
}
{ =} {  \left\langle  \partial_i \partial_j \varphi , N  \right\rangle
}
{ } {
}
{ } {
}
{ } {
}
}
{}{}{.}
Die
zweite Fundamentalmatrix
zu $\varphi$ ist die
\zusatzklammer {von

\mavergleichskettek
{\vergleichskettek
{ Q
}
{ \in }{ V
}
{  }{
}
{    }{
}
{   }{
}
}
{}{}{}} {} {}
abhängige Matrix

\mavergleichskettedisp
{\vergleichskette
{ H
}
{ =} {  \left(  h_{ij}   \right)_{1 \leq i,j \leq n-1}
}
{ } {
}
{ } {
}
{ } {
}
}
{}{}{.}
}{Die äußere Ableitung von $\omega$ wird lokal auf einer Karte, auf der $\omega$ die Gestalt

\mavergleichskettedisp
{\vergleichskette
{ \omega
}
{ =} {\sum_{ { \# \left(   I  \right) } = k , \, I \subseteq \{1 , \ldots ,  \operatorname{dim}\, M \}  }f_I dx_I
}
{ } {
}
{ } {
}
{ } {
}
}
{}{}{}
besitzt, durch

\mavergleichskettedisp
{\vergleichskette
{  d \omega
}
{ =} {\sum_{ { \# \left(   I  \right) } = k   } d(f_I ) \wedge dx_I
}
{ } {
}
{ } {
}
{ } {
}
}
{}{}{}
definiert.
}{Eine Familie von Funktionen
\maabbdisp {h_j} {X} {\R
} {}
mit

\mavergleichskette
{\vergleichskette
{ j
}
{ \in }{ J
}
{  }{
}
{    }{
}
{   }{
}
}
{}{}{}
heißt eine der Überdeckung untergeordnete Partition der Eins, wenn folgende Eigenschaften gelten.
\aufzaehlungvier{Es ist

\mavergleichskette
{\vergleichskette
{ h_j(X)
}
{ \subseteq }{ [0,1]
}
{  }{
}
{    }{
}
{   }{
}
}
{}{}{}
für alle

\mavergleichskette
{\vergleichskette
{ j
}
{ \in }{ J
}
{  }{
}
{    }{
}
{   }{
}
}
{}{}{.}
}{Jeder Punkt

\mavergleichskette
{\vergleichskette
{ P
}
{ \in }{ X
}
{  }{
}
{    }{
}
{   }{
}
}
{}{}{}
besitzt eine
\definitionsverweis {offene Umgebung}{}{}

\mavergleichskette
{\vergleichskette
{ P
}
{ \in }{  U
}
{  }{
}
{    }{
}
{   }{
}
}
{}{}{}
derart, dass die
\definitionsverweis {eingeschränkten Funktionen}{}{}

\mathl{h_j  \vert_{ U }}{} bis auf endlich viele Ausnahmen die Nullfunktion sind.
}{Es ist

\mavergleichskette
{\vergleichskette
{  \sum_{j \in J} h_j
}
{ = }{ 1
}
{  }{
}
{    }{
}
{   }{
}
}
{}{}{.}
}{Für jedes

\mavergleichskette
{\vergleichskette
{ j
}
{ \in }{ J
}
{  }{
}
{    }{
}
{   }{
}
}
{}{}{}
gibt es eine offene Menge
\mathl{W_{i(j)}}{} aus der Überdeckung derart, dass der
\definitionsverweis {Träger}{}{}
von $h_j$ in
\mathl{W_{i(j)}}{} liegt.
}
}

 }

__NOEDITSECTION__

\inputaufgabeklausurloesung
{3}

{

Formuliere die folgenden Sätze.
\aufzaehlungdrei{Der
\stichwort {Satz über die algebraische Eigenschaft der Weingartenabbildung} {.}}{Der Satz über Orientierbarkeit und positive Volumenform.}{Der
\stichwort {Brouwersche Fixpunktsatz} {.}}

}
{

\aufzaehlungdrei{\faktsituation {Es sei

\mavergleichskette
{\vergleichskette
{ W
}
{ \subseteq }{ \R^n
}
{  }{
}
{    }{
}
{   }{
}
}
{}{}{}
\definitionsverweis {offen}{}{,}
\maabb {h} { W } { \R
} {}
eine zweifach
\definitionsverweis {stetig differenzierbare Funktion}{}{}
und

\mavergleichskette
{\vergleichskette
{ Y
}
{ = }{ h^{-1}(c)
}
{  }{
}
{    }{
}
{   }{
}
}
{}{}{}
die Faser zu

\mavergleichskette
{\vergleichskette
{ c
}
{ \in }{ \R
}
{  }{
}
{    }{
}
{   }{
}
}
{}{}{,}
wobei $h$ in jedem Punkt von $Y$
\definitionsverweis {regulär}{}{}
sei.}
\faktfolgerung {Dann ist die
\definitionsverweis {Weingartenabbildung}{}{}
$L_P$ in jedem Punkt

\mavergleichskette
{\vergleichskette
{ P
}
{ \in }{ Y
}
{  }{
}
{    }{
}
{   }{
}
}
{}{}{}
\definitionsverweis {selbstadjungiert}{}{.}}
\faktzusatz {}
\faktzusatz {}}{\faktsituation {Es sei $M$ eine
\definitionsverweis {differenzierbare Mannigfaltigkeit
}{}{}
mit einer
\definitionsverweis {abzählbaren Basis der Topologie}{}{.}}
\faktfolgerung {Dann existiert genau dann eine
\definitionsverweis {stetige}{}{}
nullstellenfreie
\definitionsverweis {Volumenform}{}{}
auf $M$, wenn $M$
\definitionsverweis {orientierbar}{}{}
ist.}
\faktzusatz {Diese Volumenform kann dann auch stetig differenzierbar und positiv gewählt werden.}
\faktzusatz {}}{Es sei
\maabbdisp {\psi} { B  \left( 0,r \right) } { B  \left( 0,r \right)
} {}
eine stetig differenzierbare Abbildung der abgeschlossenen Kugel im $\R^n$ in sich. Dann besitzt $\psi$ einen Fixpunkt.}

 }

__NOEDITSECTION__

\inputaufgabeklausurloesung
{6}

{

Bestimme die
\definitionsverweis {Krümmung}{}{}
in jedem Punkt des durch

\mavergleichskettedisp
{\vergleichskette
{ h(x,y)
}
{ =} { x^2+y^2
}
{ =} { 1
}
{ } {
}
{ } {
}
}
{}{}{}
implizit gegebenen Einheitskreises mit
[[Ebene differenzierbare Kurve/Faser/Krümmung/Fakt|Lemma 3.11 (Differentialgeometrie (Osnabrück 2023))]]
und dem Gradientenfeld zu $h$.

}
{

Das
\definitionsverweis {Gradientenfeld}{}{}
zu $h$ ist $\begin{pmatrix}  2x \\ 2y   \end{pmatrix}$, ein Einheitsnormalenfeld ist somit

\mavergleichskettedisp
{\vergleichskette
{ N \begin{pmatrix}  x  \\ y   \end{pmatrix}
}
{ =} { { \frac{   1 }{  \sqrt{x^2+y^2}  } } \begin{pmatrix}  x  \\ y   \end{pmatrix}
}
{ } {
}
{ } {
}
{ } {
}
}
{}{}{.}
Das
\definitionsverweis {totale Differential}{}{}
von $N$ in einem Punkt

\mavergleichskettedisp
{\vergleichskette
{ P
}
{ =} { \begin{pmatrix}  x  \\ y   \end{pmatrix}
}
{ } {
}
{ } {
}
{ } {
}
}
{}{}{}
ist

\mavergleichskettedisp
{\vergleichskette
{  \left(D N \right)_{P}
}
{ =} { \begin{pmatrix}  { \frac{   y^2 }{   \left(  x^2+y^2  \right)^{3/2}  } }   &   - { \frac{   xy }{   \left(  x^2+y^2  \right)^{3/2}  } }   \\    - { \frac{   xy }{   \left(  x^2+y^2  \right)^{3/2}  } }    &  { \frac{   x^2 }{   \left(  x^2+y^2  \right)^{3/2}  } }  \end{pmatrix}
}
{ } {
}
{ } {
}
{ } {
}
}
{}{}{.}
Angewendet auf den tangentialen Vektor
\mathl{\begin{pmatrix}  y  \\ -x   \end{pmatrix}}{} ergibt sich

\mavergleichskettealign
{\vergleichskettealign
{ \left(D N \right)_{P} \begin{pmatrix}  y  \\ -x   \end{pmatrix}
}
{ =} { \begin{pmatrix}  y^2    &   - xy    \\  - xy   &  x^2  \end{pmatrix} \begin{pmatrix}  y  \\ -x   \end{pmatrix}
}
{ =} { \begin{pmatrix}  y \left(  y^2 +x^2  \right)  \\ -x \left(  y^2 +x^2  \right)    \end{pmatrix}
}
{ =} { \begin{pmatrix}  y   \\ -x    \end{pmatrix}
}
{ } {
}
}
{}
{}{.}
Daher ist die Krümmung der Kurve in $P$ nach
[[Ebene differenzierbare Kurve/Faser/Krümmung/Fakt|Lemma 3.11 (Differentialgeometrie (Osnabrück 2023))]]
gleich $-1$.
\zusatzklammer {Dies passt, da
\mathl{\begin{pmatrix}  y  \\ -x   \end{pmatrix}}{} und das Gradientenfeld die Standardorientierung repräsentieren und
\mathl{\begin{pmatrix}  y  \\ -x   \end{pmatrix}}{} die Tangentialrichtung für die Bewegung mit dem Uhrzeigersinn ist.} {} {}

 }

__NOEDITSECTION__

\inputaufgabeklausurloesung
{3}

{

Es sei $X$ ein
\definitionsverweis {Hausdorffraum}{}{.}
Zeige, dass die
\definitionsverweis {Diagonale}{}{}

\mavergleichskettedisp
{\vergleichskette
{ \triangle
}
{ =} { \left\{ (x,y) \in X \times X  \mid   x = y        \right\}
}
{ } {
}
{ } {
}
{ } {
}
}
{}{}{}
eine
\definitionsverweis {abgeschlossene Teilmenge}{}{}
im
\definitionsverweis {Produktraum}{}{}

\mathl{X \times X}{} ist.

}
{

Wir müssen zeigen, dass das Komplement

\mavergleichskettedisp
{\vergleichskette
{W
}
{ =} { X \times X \setminus \triangle
}
{ =} { \left\{  (x,y)   \mid   x,y \in X , \,  x \neq y       \right\}
}
{ } {
}
{ } {
}
}
{}{}{}
offen ist. Es sei also
\mathl{(x,y)}{} ein Paar mit

\mavergleichskette
{\vergleichskette
{ x
}
{ \neq }{ y
}
{  }{
}
{    }{
}
{   }{
}
}
{}{}{.}
Aufgrund der vorausgesetzten Hausdorff-Eigenschaft gibt es disjunkte offene Mengen
\mathl{U,V \subseteq X}{} mit
\mathl{x \in U}{} und
\mathl{y \in V}{.} Es ist
\mathl{(x,y) \in U \times V}{} und nach Definition der Produkttopologie ist
\mathl{U \times V}{} eine offene Menge in
\mathl{X \times X}{.} Wegen der Disjunktheit folgt aus
\mathl{(z,w) \in U\times V}{} sofort

\mavergleichskette
{\vergleichskette
{z
}
{ \neq }{w
}
{  }{
}
{    }{
}
{   }{
}
}
{}{}{.}
Also ist

\mavergleichskettedisp
{\vergleichskette
{(x,y)
}
{ \in} { U \times V
}
{ \subseteq} {  W
}
{ } {
}
{ } {
}
}
{}{}{}
und $W$ ist die Vereinigung von solchen offenen Produktmengen, also selbst offen.

 }

__NOEDITSECTION__

\inputaufgabeklausurloesung
{3}

{

Zeige, dass die drei eindimensionalen Mannigfaltigkeiten
\mathlistdisp {S^1 = \left\{ (x,y) \in \R^2  \mid   \betrag { (x,y) } = 1         \right\}} {} {\R} {und} {\R \setminus \{0\}} {}
paarweise nicht homöomorph sind.

}
{

Als abgeschlossene beschränkte Teilmenge des $\R^2$ ist der Einheitskreis kompakt
\zusatzklammer {[[Kompaktheit/Satz von Heine-Borel/Fakt|Satz von Heine-Borel]]} {} {.} Die reelle Gerade $\R$ und $\R \setminus \{0\}$ sind hingegen nicht kompakt, da sie unbeschränkte Teilmengen von $\R$ sind. Also ist
\mathl{S^1 \not \cong \R,\, \R \setminus \{0\}}{.}

Die reelle Gerade ist zusammenhängend, wie aus dem
[[Reelle Analysis/Zwischenwertsatz/Fakt|Zwischenwertsatz]]
folgt. Dagegen ist $\R \setminus \{0\}$ nicht zusammenhängend, da man

\mavergleichskettedisp
{\vergleichskette
{\R_+
}
{ =} {(\R \setminus \{0\}) \cap [0, \infty]
}
{ =} {(\R \setminus \{0\}) \cap \, ]0, \infty]
}
{ } {
}
{ } {
}
}
{}{}{}
schreiben kann, sodass man eine sowohl offene als auch abgeschlossene nichtleere Teilmenge erhält. Also ist
\mathl{\R \,\not \cong \R \setminus \{0\}}{.}

 }

__NOEDITSECTION__

\inputaufgabeklausurloesung
{5}

{

Es sei

\mavergleichskette
{\vergleichskette
{ G
}
{ \subseteq }{ \R^m
}
{  }{
}
{    }{
}
{   }{
}
}
{}{}{}
\definitionsverweis {offen}{}{}
und sei

\mavergleichskette
{\vergleichskette
{ M
}
{ \subseteq }{ G
}
{  }{
}
{    }{
}
{   }{
}
}
{}{}{}
eine
$n$-\definitionsverweis {dimensionale}{}{}
\definitionsverweis {abgeschlossene Untermannigfaltigkeit}{}{,}
die
\definitionsverweis {orientiert}{}{}
und mit der induzierten riemannschen Struktur und der
\definitionsverweis {kanonischen Volumenform}{}{}
$\omega$ versehen sei. Es sei

\mavergleichskette
{\vergleichskette
{ W
}
{ \subseteq }{ \R^n
}
{  }{
}
{    }{
}
{   }{
}
}
{}{}{}
offen und es sei
\maabbdisp {\varphi} { W } { U
} {}
ein
\definitionsverweis {Diffeomorphismus}{}{}
mit der offenen Menge

\mavergleichskette
{\vergleichskette
{ U
}
{ \subseteq }{ M
}
{  }{
}
{    }{
}
{   }{
}
}
{}{}{.}
Zeige, dass auf $W$ die Formel

\mavergleichskettealignhandlinks
{\vergleichskettealignhandlinks
{\varphi^* \left(  \omega \vert_U  \right)
}
{ =} { \left( \det  \left( \left\langle  \begin{pmatrix}  { \frac{   \partial  \varphi_1   }{  \partial   x_i   } }  \\ \vdots \\  { \frac{   \partial  \varphi_m   }{  \partial   x_i   } }   \end{pmatrix}  ,  \begin{pmatrix}  { \frac{   \partial  \varphi_1   }{  \partial   x_j   } }  \\ \vdots \\  { \frac{   \partial  \varphi_m   }{  \partial   x_j   } }   \end{pmatrix}   \right\rangle \right)_{1 \leq i,j \leq n} \right)^{1/2}  dx_1 \wedge \ldots \wedge dx_n
}
{ =} { \left( \det  \left(  { \frac{   \partial  \varphi_1   }{  \partial   x_i   } } { \frac{   \partial  \varphi_1   }{  \partial   x_j   } }  + \cdots +  { \frac{   \partial  \varphi_m   }{  \partial   x_i   } } { \frac{   \partial  \varphi_m   }{  \partial   x_j   } }  \right)_{1 \leq i,j \leq n} \right)^{1/2}  dx_1 \wedge \ldots \wedge dx_n
}
{ } {
}
{ } {
}
}
{}
{}{}
gilt.

}
{

Die zweite Gleichung ergibt sich einfach durch Auswertung des Standardskalarproduktes auf dem $\R^m$. Nach Definition der
\definitionsverweis {metrischen Fundamentalmatrix}{}{}
ist für

\mavergleichskette
{\vergleichskette
{ Q
}
{ \in }{ W
}
{  }{
}
{    }{
}
{   }{
}
}
{}{}{}

\mavergleichskettealign
{\vergleichskettealign
{ g_{ij} (Q)
}
{ =} { \left\langle  T_Q(\varphi)(e_i)  ,  T_Q(\varphi)(e_j)   \right\rangle_{\varphi(Q)}
}
{ =} { \left\langle  T_Q(\varphi)(e_i)  ,  T_Q(\varphi)(e_j)   \right\rangle
}
{ =} { \left\langle  \left(  D \varphi  \right)_{ Q } \left(   e_i   \right)  ,  \left(  D\varphi  \right)_{ Q } \left(   e_j   \right)   \right\rangle
}
{ =} { \left\langle  \begin{pmatrix}  { \frac{   \partial  \varphi_1   }{  \partial   x_i   } } ( Q)  \\ \vdots \\  { \frac{   \partial  \varphi_m   }{  \partial   x_i   } } ( Q)   \end{pmatrix}  ,  \begin{pmatrix}  { \frac{   \partial   \varphi_1   }{  \partial   x_j   } } ( Q )  \\ \vdots \\  { \frac{   \partial  \varphi_m   }{  \partial   x_j   } } ( Q)   \end{pmatrix}   \right\rangle
}
}
{}
{}{,}
da ja der Tangentialraum
\mathl{T_{\varphi(Q)}M}{} das induzierte Skalarprodukt des $\R^m$ trägt, da die Tangentialabbildung im lokalen Fall das totale Differential ist und da man dessen Einträge mit den partiellen Ableitungen ausdrücken kann. Daher ergibt sich die Behauptung aus
[[Riemannsche Mannigfaltigkeit/Orientiert/Kanonische Volumenform/Lokale Berechnung/Fakt|Lemma 16.7 (Differentialgeometrie (Osnabrück 2023))]].

 }

__NOEDITSECTION__

\inputaufgabeklausurloesung
{weiter}

{

Wir betrachten den

\mavergleichskette
{\vergleichskette
{ \R^6
}
{ = }{ \left(  \R^2  \right)^3
}
{  }{
}
{    }{
}
{   }{
}
}
{}{}{}
als Menge aller
\zusatzklammer {auch entarteter} {} {}
Dreiecke, indem wir ein Dreieck mit den
\zusatzklammer {geordneten} {} {}
Eckpunkten

\mavergleichskette
{\vergleichskette
{ P_1
}
{ = }{ \begin{pmatrix}  x_1  \\ y_1    \end{pmatrix}
}
{  }{
}
{    }{
}
{   }{
}
}
{}{}{,}

\mavergleichskette
{\vergleichskette
{ P_2
}
{ = }{ \begin{pmatrix}  x_2  \\ y_2    \end{pmatrix}
}
{  }{
}
{    }{
}
{   }{
}
}
{}{}{}
und

\mavergleichskette
{\vergleichskette
{ P_3
}
{ = }{ \begin{pmatrix}  x_3  \\ y_3    \end{pmatrix}
}
{  }{
}
{    }{
}
{   }{
}
}
{}{}{,}
mit dem Koordinatentupel

\mathdisp {(x_1,y_1,x_2,y_2,x_3,y_3)} {  }

identifizieren.
\aufzaehlungsechs{Zeige, dass die Menge der Dreiecke, bei denen zwei Eckpunkte zusammenfallen, eine abgeschlossene Teilmenge des $\R^6$ ist
\zusatzklammer {das Komplement davon ist somit eine offene Menge in $\R^6$, die wir $U$ nennen} {} {.}
}{Zeige, dass die Menge der Dreiecke, bei denen alle drei Eckpunkte auf einer Geraden liegen, eine abgeschlossene Teilmenge des $\R^6$ ist
\zusatzklammer {das Komplement davon, das aus allen nichtentarteten Dreiecken besteht, ist somit eine offene Menge in $\R^6$, die wir $V$ nennen} {} {.}
}{Erstelle eine Funktionsvorschrift, die die Abbildung
\maabbdisp {F} {\R^6} {\R
} {}
beschreibt, die einem Dreieck seinen Umfang zuordnet.
}{Zeige, dass die Funktion $F$ aus Teil (3) auf der Menge $U$ stetig differenzierbar ist.
}{Berechne die partielle Ableitung von $F$ nach $x_1$ auf $U$.
}{Zeige, dass die Menge der nichtentarteten Dreiecke mit Umfang $1$ eine
\definitionsverweis {abgeschlossene Untermannigfaltigkeit}{}{}
von $V$ bildet. Was ist die Dimension?
}

}
{

\aufzaehlungsechs{Die Eckpunkte
\mathkor {} {P_1} {und} {P_2} {}
stimmen genau dann überein, wenn ihre beiden Koordinaten übereinstimmen, wenn also

\mavergleichskette
{\vergleichskette
{ x_1
}
{ = }{ x_2
}
{  }{
}
{    }{
}
{   }{
}
}
{}{}{}
und

\mavergleichskette
{\vergleichskette
{ y_1
}
{ = }{ y_2
}
{  }{
}
{    }{
}
{   }{
}
}
{}{}{}
ist. Eine solche Bedingung definiert eine abgeschlossene Teilmenge, da es sich um die Faser einer
\zusatzklammer {stetigen} {} {}
linearen Abbildung handelt. Da eine endliche Vereinigung von abgeschlossenen Teilmengen wieder abgeschlossen ist, ist die in Frage stehende Menge abgeschlossen.
}{Wir betrachten die Verbindungsvektoren

\mathdisp {P_2-P_1 = \begin{pmatrix}  x_2-x_1  \\ y_2-y_1   \end{pmatrix} \text{   und } P_3-P_1 = \begin{pmatrix}  x_3-x_1  \\ y_3-y_1   \end{pmatrix}} {  }

Die drei Punkte sind genau dann kollinear, wenn diese zwei Vektoren linear abhängig sind, und dies ist genau dann der Fall, wenn ihre Determinante $0$ ist. Die Kollinearitätssbedingung lautet somit

\mavergleichskettedisp
{\vergleichskette
{ H(x_1,x_2,x_3,y_1,y_2,y_3)
}
{ =} { (x_2-x_1) (y_3-y_1) - (x_3-x_1) (y_2-y_1)
}
{ =} { 0
}
{ } {
}
{ } {
}
}
{}{}{.}
Da $H$ eine polynomiale Abbildung und damit stetig ist, ist die Faser zu $0$ eine abgeschlossene Teilmenge.
}{Da der Umfang einfach die Summe der drei beteiligten Dreiecksseiten ist, gilt

\mavergleichskettealigndrucklinks
{\vergleichskettealigndrucklinks
{ F(x_1,x_2,x_3,y_1,y_2,y_3)
}
{ =} { d(P_1,P_2) + d(P_1,P_3) +d(P_2,P_3)
}
{ =} { \sqrt{(x_2-x_1)^2 + (y_2-y_1)^2} +  \sqrt{(x_3-x_1)^2 + (y_3-y_1)^2} +  \sqrt{(x_3-x_2)^2 + (y_3-y_2)^2}
}
{ } {
}
{ } {
}
}
{}{}{.}
}{Für einen Punkt aus
\mathl{U}{} ist jeder Radikand als polynomiale Funktion stetig differenzierbar und echt positiv, somit sind auch die Quadratwurzeln daraus stetig differenzierbar.
}{Die partielle Ableitung ist

\mavergleichskettealigndrucklinks
{\vergleichskettealigndrucklinks
{ { \frac{   \partial   F   }{  \partial   x_1  } }
}
{ =} { { \frac{   1 }{  2 } }   \left(  (  x_2-x_1  )^2 + (y_2-y_1)^2  \right)^{-1/2} (2x_1-2x_2)    + { \frac{   1 }{  2 } }   \left(  (  x_3-x_1  )^2 + (y_3-y_1)^2  \right)^{-1/2} (2x_1-2x_3)
}
{ =} {  \left(  (  x_2-x_1  )^2 + (y_2-y_1)^2  \right)^{-1/2} (x_1-x_2)    +   \left(  (  x_3-x_1  )^2 + (y_3-y_1)^2  \right)^{-1/2} (x_1-x_3)
}
{ } {
}
{ } {
}
}
{}{}{.}
}{Auf $U$ sind die Radikanden positiv, somit verschwindet die partielle Ableitung nach $x_1$ nur, falls
\mathl{x_1=x_2=x_3}{} ist. Ein Dreieck mit dieser Bedingung ist aber kollinear. Das bedeutet, dass auf $V$ die partielle Ableitung nach $x_1$ nie $0$ wird und das bedeutet insbesondere, dass der Gradient auf $V$ nirgends verschwindet, also regulär ist. Nach
[[Satz über implizite Abbildungen/Faser ist Mannigfaltigkeit/Fakt|Satz 8.2 (Differentialgeometrie (Osnabrück 2023))]]
ist die Faser eine abgeschlossene Untermannigfaltigkeit, und zwar der Dimension

\mavergleichskettedisp
{\vergleichskette
{ 6-1
}
{ =} { 5
}
{ } {
}
{ } {
}
{ } {
}
}
{}{}{.}
}

 }

__NOEDITSECTION__

\inputaufgabeklausurloesung
{4}

{

Es seien
\mathkor {} {r} {und} {R} {}
positive reelle Zahlen mit

\mavergleichskette
{\vergleichskette
{ 0
}
{ < }{ r
}
{ < }{ R
}
{    }{
}
{   }{
}
}
{}{}{,}
wir betrachten den
\zusatzklammer {eingebetteten} {} {}
Torus

\mavergleichskettedisp
{\vergleichskette
{ T
}
{ =} { \left\{ (x,y,z) \in \R^3  \mid   \left(  \sqrt{x^2+y^2} -R  \right)^2 +z^2  = r^2         \right\}
}
{ \subseteq} { \R^3
}
{ } {
}
{ } {
}
}
{}{}{}
mit der
\definitionsverweis {induzierten riemannschen Struktur}{}{.}
Zeige, dass zu jedem

\mavergleichskette
{\vergleichskette
{ \alpha
}
{ \in }{ [0,2 \pi [
}
{  }{
}
{    }{
}
{   }{
}
}
{}{}{}
die Abbildung
\maabbeledisp {\Psi_\alpha} {\R^3} {\R^3
} { \left(  x   , \,   y  , \,  z                 \right) } { \left(  x \cos \alpha - y \sin  \alpha   , \,   x \sin \alpha + y \cos  \alpha  , \,  z                 \right)
} {,}
eine
\definitionsverweis {Isometrie}{}{}
von $T$ in sich induziert.

}
{

Die Abbildung ist auf

\mavergleichskettedisp
{\vergleichskette
{ \R^3
}
{ =} { \R^2 \times \R
}
{ } {
}
{ } {
}
{ } {
}
}
{}{}{}
das Produkt einer
\definitionsverweis {ebenen Drehung}{}{}
mit der Identität. Daher liegt eine euklidische
\definitionsverweis {Isometrie}{}{}
des $\R^3$ vor. Es sei

\mavergleichskette
{\vergleichskette
{ (x,y,z)
}
{ \in }{  T
}
{  }{
}
{    }{
}
{   }{
}
}
{}{}{}
und

\mavergleichskettedisp
{\vergleichskette
{ \Psi_\alpha(x,y,z)
}
{ =} { ( u,v,z)
}
{ } {
}
{ } {
}
{ } {
}
}
{}{}{.}
Dann ist

\mavergleichskettealigndrucklinks
{\vergleichskettealigndrucklinks
{ \left(  \sqrt{u^2+v^2} -R  \right)^2 +z^2
}
{ =} { \left(  \sqrt{ \left(  x \cos \alpha - y \sin  \alpha   \right)^2+ \left(  x \sin \alpha + y \cos  \alpha  \right)^2} -R  \right)^2 +z^2
}
{ =} { \left(  \sqrt{  x^2 \cos^{ 2  } \alpha + y^2 \sin^{ 2  }  \alpha  -2xy \cos \alpha \sin  \alpha    + x^2 \sin^{ 2  } \alpha + y^2 \cos^{ 2  }  \alpha  +2xy \cos \alpha \sin  \alpha     } -R  \right)^2 +z^2
}
{ =} { \left(  \sqrt{  x^2   + y^2 } -R  \right)^2 +z^2
}
{ =} { r^2
}
}
{}{}{,}
also

\mavergleichskette
{\vergleichskette
{ \Psi_\alpha(x,y,z)
}
{ \in }{  T
}
{  }{
}
{    }{
}
{   }{
}
}
{}{}{.}
Es ist also
\maabbdisp {\Psi_\alpha} { T} {T
} {}
ein Diffeomorphismus und eine Isometrie, da $T$ die induzierte riemannsche Struktur trägt.

 }

__NOEDITSECTION__

\inputaufgabeklausurloesung
{3}

{

Beschreibe ein Modell für die hyperbolische Fläche.

}
{  }

__NOEDITSECTION__

\inputaufgabeklausurloesung
{7}

{

Beweise das Theorema egregium.

}
{

Wir differenzieren die erste Bestimmungsgleichung für die Christoffelsymbole, also

\mavergleichskettedisp
{\vergleichskette
{ \partial_1 \partial_1 \varphi
}
{ =} { \Gamma^1_{11}  \partial_1 \varphi  + \Gamma^2_{11}  \partial_2 \varphi +h_{11} N
}
{ } {
}
{ } {
}
{ } {
}
}
{}{}{,}
in Richtung der zweiten Variablen und die zweite Bestimmungsgleichung, also

\mavergleichskettedisp
{\vergleichskette
{ \partial_1 \partial_2 \varphi
}
{ =} { \Gamma^1_{12}  \partial_1 \varphi + \Gamma^2_{12}  \partial_2  \varphi + h_{12} N
}
{ } {
}
{ } {
}
{ } {
}
}
{}{}{,}
in Richtung der ersten Variablen und erhalten nach Schwarz

\mavergleichskettealigndrucklinks
{\vergleichskettealigndrucklinks
{ \left(  \partial_2 \Gamma^1_{11}  \right) \partial_1 \varphi +  \left(  \partial_2 \Gamma^2_{11}  \right)  \partial_2 \varphi+ \left(  \partial_2 h_{11}  \right)  N + \Gamma^1_{11} \partial_2 \partial_1 \varphi + \Gamma^2_{11}  \partial_2 \partial_2 \varphi  +h_{11} \partial_2 N
}
{ =} { \left(  \partial_1 \Gamma^1_{12}  \right)  \partial_1  \varphi+ \left(  \partial_1 \Gamma^2_{12}  \right)\partial_2 \varphi + \left(  \partial_1 h_{12}  \right)  N+ \Gamma^1_{12}  \partial_1 \partial_1 \varphi + \Gamma^2_{12}  \partial_1 \partial_2 \varphi + h_{12} \partial_1 N
}
{ } {
}
{ } {
}
{ } {
}
}
{}{}{.}
Die Differenz dieser Ausdrücke ist $0$, und wir bestimmen, was sich dabei auf den Basisvektor $\partial_2 \varphi$ bezieht. Dazu müssen wir die Bestimmungsgleichungen für die Christoffelsymbole und
[[Hyperfläche/Raum/Parametrisierung/Weingartenabbildung/Fundamentalformen/Fakt|Lemma 19.3 (Differentialgeometrie (Osnabrück 2023))  (3)]]
heranziehen und erhalten

\mavergleichskettedisphandlinks
{\vergleichskettedisphandlinks
{ 0
}
{ =} { \partial_2 \Gamma^2_{11} - \partial_1 \Gamma^2_{12} +\Gamma^1_{11} \Gamma^2_{12} - \Gamma^1_{12} \Gamma^2_{11} + \Gamma_{11}^2 \Gamma_{22}^2- \left(  \Gamma_{12}^2  \right)^2 + h_{11}  { \frac{   g_{12} h_{12} - g_{11} h_{22}   }{  g_{11}g_{22} -g_{12}^2 } } - h_{12}  { \frac{   g_{12} h_{11} - g_{11}h_{12}  }{  g_{11}g_{22} -g_{12}^2 } }
}
{ } {
}
{ } {
}
{ } {
}
}
{}{}{.}
Mit

\mavergleichskettedisphandlinks
{\vergleichskettedisphandlinks
{ - g_{11} \left(  h_{11}h_{22} -h_{12}^2  \right)
}
{ =} { h_{11} \left(  g_{12} h_{12}- g_{11} h_{22}  \right) - h_{12} \left(  g_{12} h_{11} - g_{11}h_{12}  \right)
}
{ } {
}
{ } {
}
{ } {
}
}
{}{}{}
können wir die beiden hinteren Summanden ersetzen und erhalten mit
[[Hyperfläche/Raum/Parametrisierung/Weingartenabbildung/Fundamentalformen/Fakt|Lemma 19.3 (Differentialgeometrie (Osnabrück 2023))  (4)]]

\mavergleichskettealignhandlinks
{\vergleichskettealignhandlinks
{ g_{11} K
}
{ =} { g_{11} { \frac{   h_{11}h_{22} -h_{12}^2  }{ g_{11}g_{22} -g_{12}^2 } }
}
{ =} { \left(  \partial_2 \Gamma_{11}^2 - \partial_1 \Gamma_{12}^2 + \Gamma^1_{11} \Gamma^2_{12}  + \Gamma_{11}^2 \Gamma_{22}^2 - \Gamma_{12}^1 \Gamma_{11}^2-  \left(  \Gamma_{12}^2  \right)^2  \right)
}
{ } {
}
{ } {
}
}
{}
{}{.}

 }

__NOEDITSECTION__

\inputaufgabeklausurloesung
{6 (1+2+2+1)}

{

Wir betrachten die
\definitionsverweis {Differentialform}{}{}

\mavergleichskettedisp
{\vergleichskette
{ \omega
}
{ =} { ydx+zdy +xdz
}
{ } {
}
{ } {
}
{ } {
}
}
{}{}{}
auf dem $\R^3$ und die Abbildung
\maabbeledisp {\varphi} {\R^2} {\R^3
} { \left(  u   , \,  v                   \right) } { \left( u^2  , \,  v^2  , \,  uv                 \right)
} {.}
\aufzaehlungvier{Berechne die äußere Ableitung von $\omega$.
}{Berechne den Rückzug $\varphi^* \omega$  von $\omega$ unter $\varphi$.
}{Berechne die äußere Ableitung von $\varphi^* \omega$ auf $\R^2$.
}{Berechne den Rückzug $\varphi^* d \omega$  von $d\omega$ unter $\varphi$ unabhängig von (3).
}

}
{

\aufzaehlungvier{Es ist

\mavergleichskettedisp
{\vergleichskette
{ d \omega
}
{ =} { d \left(  ydx+zdy +xdz  \right)
}
{ =} { dy \wedge dx +dz \wedge dy +dx \wedge dz
}
{ =} { - dx \wedge dy -dy \wedge dz +dx \wedge dz
}
{ } {
}
}
{}{}{.}
}{Es ist

\mavergleichskettealign
{\vergleichskettealign
{  \varphi^* \omega
}
{ =} { v^2 d u^2 + uvdv^2+ u^2duv
}
{ =} { 2v^2udu +2uv^2 dv + u^2(udv +vdu)
}
{ =} { \left(  2uv^2 +u^2v  \right) du + \left(  2uv^2 +u^3  \right) dv
}
{ } {
}
}
{}
{}{.}
}{Es ist

\mavergleichskettealign
{\vergleichskettealign
{  d \left(   \left(  2uv^2 +u^2v  \right) du + \left(  2uv^2 +u^3  \right) dv     \right)
}
{ =} {  d \left(  2uv^2 +u^2v  \right) du  +  d\left(  2uv^2 +u^3  \right) dv
}
{ =} { \left(  4uv  + u^2      \right)  dv \wedge du  +   \left(  2v^2   +3u^2  \right)   du \wedge dv
}
{ =} { \left(  -4uv +  2v^2   +2u^2  \right)   du \wedge dv
}
{ } {
}
}
{}
{}{.}
}{Der Rückzug $\varphi^* d \omega$ ist

\mavergleichskettealign
{\vergleichskettealign
{ \varphi^* \left(  - dx \wedge dy -dy \wedge dz +dx \wedge dz   \right)
}
{ =} { - d u^2 \wedge dv^2 -d v^2 \wedge duv +d u^2 \wedge duv
}
{ =} { -4uv  du \wedge dv -2v^2 dv \wedge du + 2   u^2 du \wedge dv
}
{ =} { \left(  -4uv+2v^2 + 2 u^2  \right)  du \wedge dv
}
{ } {
}
}
{}
{}{.}
}

 }

__NOEDITSECTION__

\inputaufgabeklausurloesung
{3}

{

Es sei $M$ eine
\definitionsverweis {differenzierbare Mannigfaltigkeit}{}{}
und $\omega$ eine
\definitionsverweis {exakte}{}{} \definitionsverweis {Differentialform}{}{}
auf $M$. Es sei $S$ eine
\definitionsverweis {kompakte}{}{}
\definitionsverweis {orientierte}{}{}
\definitionsverweis {differenzierbare Mannigfaltigkeit}{}{}
\zusatzklammer {ohne Rand} {} {}
mit
\definitionsverweis {abzählbarer Basis der Topologie}{}{}
und es sei
\maabbdisp {\varphi} { S } { M
} {}
eine
\definitionsverweis {stetig differenzierbare}{}{}
Abbildung. Zeige

\mavergleichskettedisp
{\vergleichskette
{
\int_{  S  }   \varphi^*\omega
}
{ =} { 0
}
{ } {
}
{ } {
}
{ } {
}
}
{}{}{.}

}
{

Da nach
[[Differentialform/Mannigfaltigkeit/Äußere Ableitung/Eigenschaften/Fakt|Satz 20.4 (Differentialgeometrie (Osnabrück 2023))  (5)]]
das Zurückziehen von Formen mit der äußeren Ableitung verträglich ist, ist auch
\mathl{\varphi^*\omega}{} exakt. Es sei $\tau$ eine differenzierbare Form auf $S$ mit

\mavergleichskettedisp
{\vergleichskette
{d \tau
}
{ =} { \varphi^* \omega
}
{ } {
}
{ } {
}
{ } {
}
}
{}{}{.}
Nach
dem [[Mannigfaltigkeit mit Rand/Satz von Stokes/Fakt|Satz von Stokes]]
ist

\mavergleichskettedisp
{\vergleichskette
{ \int_S \varphi^* \omega
}
{ =} { \int_S d \tau
}
{ =} { \int_{ \partial S } \tau
}
{ =} { \int_\emptyset \tau
}
{ =} { 0
}
}
{}{}{.}

 }

__NOEDITSECTION__

\inputaufgabeklausurloesung
{6}

{

Beweise den Satz über die Retraktionen auf den Rand auf Mannigfaltigkeiten.

}
{

Der Rand
\mathl{\partial M}{} ist nach
[[Mannigfaltigkeit mit Rand/Orientierung/Randorientierung/Fakt|Satz 21.8 (Differentialgeometrie (Osnabrück 2023))]]
eine
\definitionsverweis {orientierte}{}{}
\definitionsverweis {differenzierbare Mannigfaltigkeit}{}{}
\zusatzklammer {ohne Rand} {} {.}
Daher gibt es nach
[[Nullstellenfreie Volumenform/Orientierte (Karten) Mannigfaltigkeit/Äquivalenz/Fakt|Satz 22.11 (Differentialgeometrie (Osnabrück 2023))]]
eine
\definitionsverweis {stetig differenzierbare}{}{}
\definitionsverweis {positive Volumenform}{}{}
$\tau$ auf
\mathl{\partial M}{.} Es ist
\mathl{\int_{  \partial M  }  \tau >0}{.} Die
\definitionsverweis {äußere Ableitung}{}{}
der Volumenform $\tau$ ist $0$. Nehmen wir an, dass es eine
\definitionsverweis {stetig differenzierbare Abbildung}{}{}
\maabbdisp {\varphi} { M } { \partial M
} {}
mit
\mathl{\varphi \vert_{\partial M} = \operatorname{Id}_{\partial M}}{} gebe. Dann ist die
\definitionsverweis {zurückgezogene Form}{}{}

\mathl{\varphi^* \tau}{} eine
$(n-1)$-\definitionsverweis {Differentialform}{}{}
auf $M$, deren Einschränkung auf den Rand mit $\tau$ übereinstimmt. Daher gilt unter Verwendung von
[[Mannigfaltigkeit mit Rand/Satz von Stokes/Fakt|Satz 23.2 (Differentialgeometrie (Osnabrück 2023))]]
und
[[Differentialform/Mannigfaltigkeit/Äußere Ableitung/Eigenschaften/Fakt|Satz 20.4 (Differentialgeometrie (Osnabrück 2023))  (5)]]

\mavergleichskettealign
{\vergleichskettealign
{
\int_{  \partial M  }  \tau
}
{ =} {
\int_{  \partial  M  }  \varphi^* \tau
}
{ =} {
\int_{   M  }  d(\varphi^* \tau)
}
{ =} {
\int_{   M  }  \varphi^* (d\tau)
}
{ =} {
\int_{   M  }  \varphi^* (0)
}
}
{
\vergleichskettefortsetzungalign
{ =} { 0
}
{ } {}
{ } {}
{ } {}
}
{}{.}
Dies ist ein Widerspruch.

 }

 [[Kategorie:Latexseite]]
